\section{Extended Profile and Robustness Results}
\label{app:extended_results}

This appendix reports extended parameter sweeps, profile analyses, and robustness diagnostics that support the main-text results.

\subsection{Precision-dynamics perturbations}
\label{sec:perturbation_dynamics}

We first perturb the amplitude and stability of $\beta_k$ dynamics to test how precision dynamics affect action commitment. These variants are computational controls rather than diagnostic categories.

\begin{table}[H]
\centering
\caption{Precision-dynamics perturbations in the graded decision regime.}
\label{tab:perturbation_dynamics}
\begin{tabular}{lccc}
\toprule
Variant & $\beta_k$ range & Total payoff & Policy entropy \\
\midrule
affect & 1.32 & 1851.3 & 8.6 \\
low-gain & 0.24 & 1840.5 & 8.9 \\
high-gain & 1.48 & 1859.0 & 8.1 \\
cautious-prior & 1.48 & 1877.7 & 8.4 \\
\bottomrule
\end{tabular}
\end{table}

\begin{sloppypar}
Flattening the precision signal produces the clearest behavioral change. The low-gain variant compresses $\beta_k$ range to $0.24$ and raises entropy to $8.9$, near the no-affect and tracked-only baselines. Expanded-dynamics variants produce wider $\beta_k$ movement, but payoff does not follow range monotonically. The parameter sweeps below therefore vary gain and prior structure directly.
\end{sloppypar}

\begin{figure}[H]
\centering
\includegraphics[width=\linewidth]{fig_phenotype_dynamics_summary.pdf}
\caption{Precision-dynamics perturbation profiles in the graded decision regime. Variants compare precision-dynamics settings in the graded decision regime.}
\label{fig:phenotype_dynamics_appendix}
\end{figure}

\subsection{Alpha sweep}
\label{sec:phenotypes_alpha}

Sweeping $\alpha$ across $\{0.05, 0.1, 0.3, 0.5, 1.0, 2.0, 4.0, 8.0\}$ tests whether precision gain controls reactive social confidence across metrics. $\beta_k$ movement range increases monotonically with $\alpha$ in the betrayal regime:

\begin{table}[H]
\centering
\caption{$\beta_k$ range by precision-gain $\alpha$ (betrayal regime, 20 seeds).}
\label{tab:alpha_sweep}
\begin{tabular}{cc}
\toprule
$\alpha$ & Mean $\beta_k$ range \\
\midrule
0.05 & 0.147 \\
0.1  & 0.182 \\
0.3  & 0.326 \\
0.5  & 0.416 \\
1.0  & 0.580 \\
2.0  & 0.862 \\
4.0  & 1.051 \\
8.0  & 1.218 \\
\bottomrule
\end{tabular}
\end{table}

\begin{figure}[H]
\centering
\includegraphics[width=\linewidth]{fig_alpha_sweep.pdf}
\caption{Precision-gain sweep across open graded and betrayal environments (20 seeds). Higher $\alpha$ expands $\beta_k$ range and reshapes engagement readouts without monotonic payoff ranking.}
\label{fig:alpha_sweep_appendix}
\end{figure}

\subsection{Profile quadrant profiles}
\label{sec:phenotypes_quadrants}

Crossing prior belief (naive versus cautious) with $\alpha$ (low versus high) defines four behavioral profiles tested in the betrayal environment (20 seeds). Table~\ref{tab:phenotypes} gives key metrics.

\begin{table}[H]
\centering
\caption{Behavioral metrics by profile (betrayal environment, 20 seeds). Default agent shown for reference.}
\label{tab:phenotypes}
\begin{tabular}{lccccc}
\toprule
Profile & Payoff & Recovery & Gini & $\beta$ range & Trust asymmetry \\
\midrule
Anxious-reactive   & 1932 & 26.5 & 0.338 & 1.084 & 0.98 \\
Hypervigilant      & 1951 & 20.1 & 0.423 & 0.892 & 1.65 \\
Naive-stubborn     & 1894 & 21.3 & 0.304 & 0.364 & 1.19 \\
Avoidant-rigid     & 1889 & 14.9 & 0.357 & 0.322 & 1.21 \\
Default            & 1991 & 23.0 & 0.423 & 0.893 & 1.68 \\
\bottomrule
\end{tabular}
\end{table}

\begin{figure}[H]
\centering
\includegraphics[width=\linewidth]{fig_phenotype_quadrants.pdf}
\caption{Gain-prior profile quadrants from the prior-by-$\alpha$ factorial (betrayal environment, 20 seeds).}
\label{fig:phenotype_quadrants_appendix}
\end{figure}

\subsection{Forgiveness and trust repair}

The forgiveness paradigm uses a 200-round partner-choice episode: cooperative baseline (rounds 1--80), betrayal on P0 (rounds 81--120), then reversion to cooperative/trusting behavior (rounds 121--200). Readouts include reengagement latency, post-return selection, payoff recovery, and $\beta_k$ recovery trajectories for the reverted partner.

\begin{table}[H]
\centering
\caption{Forgiveness and trust-repair metrics by precision profile (20 seeds). Reengagement is the post-repair P0 selection rate; payoff recovery compares late post-repair payoff with the late pre-betrayal baseline.}
\label{tab:forgiveness}
\begin{tabular}{lcccc}
\toprule
Variant & Reengagement & Payoff recovery & Latency & $\beta_k$ range \\
\midrule
cautious-high-$\alpha$ & 0.518 & 0.996 & 5.5 & 0.971 \\
cautious-low-$\alpha$ & 0.630 & 1.014 & 1.5 & 0.287 \\
default & 0.475 & 1.019 & 3.0 & 1.012 \\
naive-high-$\alpha$ & 0.560 & 1.005 & 2.2 & 1.100 \\
naive-low-$\alpha$ & 0.527 & 1.004 & 4.6 & 0.325 \\
no-affect & 0.593 & 1.033 & 1.7 & -- \\
\bottomrule
\end{tabular}
\end{table}

\begin{figure}[H]
\centering
\includegraphics[width=\linewidth]{fig_forgiveness.pdf}
\caption{Forgiveness/trust-repair diagnostic from the 20-seed profile analyses. Reengagement and payoff recovery do not collapse into a single monotonic gain effect: no-affect and cautious-low-$\alpha$ profiles reengage most often, while payoff recovery remains near baseline across profiles.}
\label{fig:forgiveness_appendix}
\end{figure}

\subsection{Mixed-volatility extended diagnostics}

The primary mixed-volatility table and figure remain in Section~\ref{sec:phenotypes_mixed}. Extended Exp D diagnostics include per-partner $\beta_k$ trajectories, rolling P0 selection concentration, and pairwise stability--precision correlations used to compute the discrimination index and P0 false-positive rate reported in the main text.
